% Imporditud: tools/import_eksperimendid.py
% Allikas: Eesti füüsikaolümpiaadi eksperimendi ülesannete
%   kogu (Rael Kalda, 2023), ülesanne Ü147, lahendus L147.
\setAuthor{Jaan Kalda}
\setRound{piirkonnavoor}
\setYear{2016}
\setNumber{G E2}
\setDifficulty{4}
\setTopic{dünaamika}
\setVahendid{kirjaklambrid (15 tk), silinder, niit (niidi mass on tühine võrreldes kir- jaklambri massiga)}
\setPunktid{12}
\setAllikas{Eksperimendi kogu Ü147, lahendus lk 109}
\setLahendusseis{toores}

\prob{Niit}
Määrake niidi ja silindri välispinna vaheline liughõõrdetegur $\mu$ täpsusega kuni 10\%.
\textit{Vihje:} on teada, et kui ümber pulga on keeratud niit teatud arv keerde nii, et niidi sihi summaarne pöördenurk radiaanides on $\varphi$ (nt $\varphi$ = 2$\pi$ vastab täpselt ühele täiskeerule) ja niidi ühest otsast tõmmatakse jõuga $T_1$ ning teisest otsast jõuga $T_2$ < $T_1$, siis niit hakkab libisema kui $T_1$ > T2e$\mu\varphi$.

\hint

\solu
Seome niidi ühte otsa n klambrit (klambritest on mugav teha ahel) ja teise otsa m ning viskame niidi üle silindri (nii et niidi siht pöördub nurga $\pi$ võrra). Alustuseks n = 1 ja leiame sellise vähima m-i väärtuse $m_1$, mille puhul niit hakkab libisema raskema otsa poole. Teame nüüd, et hõõrdetegur rahuldab võrratusi

$m_1$ $-$ 1 < e$\mu\pi$ < $m_1$,

millest 1 ln($m_1$ $-$ 1) < $\mu$ < 1 ln($m_1$). Võtame nüüd n = 2 ning analoogselt leiame $\pi$ $\pi$ m$^2$ väärtuse; kui eelmine osa oli õigesti tehtud, siis on meil kaks võimalust: m$^2$ =

2m1 või m$^2$ = 2m1 $-$ 1. Nüüd saame võrratuseks

( ) 1 ln ( m$^2$ ) . 1 ln m$^2$ $-$ 1 < $\mu$ < $\pi$ 2 $\pi_2$

jSuõultruevsa; lütlhdõjuõhrduelt$\pi$e1glunr(i mvännä$-$r1t)us<es$\mu$t võib olnll(amvnõni)m. Saelilklekpaõtshejalälbsiavaimiaekvaanst=us3ekjas n=4 anda < 1 $\pi$

1 ( mn(mn $-$ 1) ) $\mu \approx$ 2$\pi$ ln $n_2$ .

Samuti, kui hõõrdetegur on piisavalt väike, saab teha niidiga poole keeru asemel

poolteist keerdu (pöördenurk 3$\pi$). Olgu sellisel juhul niidi libisemise ajal ühes ot-

sas 1 klamber ja teises M klambrit; siis 1 ln(M $-$ 1) < $\mu$ < 1 ln M ning 3$\pi$ 3$\pi$

$\mu \approx$ 6$\pi$ ln[(M $-$ 1)M ].

Libeda silindri ja niidi puhul võivad tulla näiteks järgmised tulemused: $m_1$ = 3, m$^2$ = 5, m$^3$ = 7 ja $m_4$ = 9. Kõige täpsema tulemuse annab n = 4, $\mu \approx$ ln(4,5)/(2$\pi$) $\approx$ 0,24. Kui viia katse läbi poolteise keeru juures, siis saame M = 11, mis annab $\mu \approx$ ln(110)/(6$\pi$) $\approx$ 0,25, mis on natuke täpsem tulemus (sest antud juhul M > $m_4$).

Hindamine: Katse idee: 3 p. Mõõtmiste eest alljärgnevalt: leitud katseliselt ainult $m_1$: 2 p; leitud katseliselt m$^2$, kuid mitte mn, kus n > 2 ($m_1$ leidmine siinjuures punkte ei lisa, sest täpsust ei suurenda): 4 p; leitud katseliselt m$^3$, kuid mitte mn, kus n > 3: 5 p; leitud katseliselt $m_4$: 6 p (kui hõõrdetegur on nii suur, et n = 4 puhul ei jagu klambreid, siis tuleb anda n = 3 eest 5 p asemel 6 p); katse läbi viidud ka poolteise keeru juures: täiendavalt 1 p. Kui hõõrdetegur on nii suur, et klambreid ei jätku katse läbiviimiseks poolteise keeru juures, siis kantakse see punkt eelmistele mõõtmistele, st n = 1: 3 p; n = 2: 5 p; n = 3: 7 p. Kui katse on läbi viidud ainult poolteise keeru puhul, kuid ei ole mõõdetud poole keeru juures, siis antakse mõõtmiste eest kokku 6 p. Tulemuste täpsuse eest: Vastus arvutatud mõõtmistulemuste põhjal vigadeta ja tulemuse erinevus tegelikust väärtusest ei ole suurem kui 10\%: 2 p; on 15\% ja 10\% vahel: 1,5 p; on 20\% ja 15\% vahel: 1 p; on 30\% ja 20\% vahel: 0,5 p.

Märkus: põhimõtteliselt on vale katset läbi viia nii, et riputada niidi otstesse võrdsed arvud klambreid, asetada see silindrile ning kallutada silindri telge leides nurga, mille juures niit hakkab libisema. Esiteks, sellisel juhul ei kehti seos $\mu$ = tan $\alpha$, teiseks, nurga mõõtmiseks pole katsevahendeid. Selline lahendus annab 0 p.
\probend
